\section*{附加习题：网}\label{section:ch03-extra-nets}

我们已经看到，要区别度量空间中的极限点、连续函数以及紧致空间，只要用序列就``足够''了。但是，对于任意的拓扑空间要做同样的事就需要一种称为\textbf{网}的概念，它是序列概念的推广。在这里我们给出有关的定义，而将证明留作练习。首先回忆一下集合 $A$ 中的一个关系 $\preceq$，如果满足下列条件，则称为偏序关系：
\begin{enumerate}
\item 对于所有 $\alpha$，$\alpha\preceq\alpha$.
\item 若 $\alpha\preceq\beta$ 和 $\beta\preceq\alpha$，则 $\alpha=\beta$.
\item 若 $\alpha\preceq\beta$ 和 $\beta\preceq\gamma$，则 $\alpha\preceq\gamma$.
\end{enumerate}

现在我们给出以下定义：

$J$ 是一个具有偏序 $\preceq$ 的集合，若满足下列条件，则称 $J$ 是\textbf{有向集}（directed set）：对于 $J$ 中每一对元素 $\alpha$ 和 $\beta$，存在 $J$ 中的一个元素 $\gamma$，具有性质 $\alpha\preceq\gamma$ 和 $\beta\preceq\gamma$。

\begin{enumerate}
\item 证明下列集合为有向集：
\begin{enumerate}
\item 在关系 $\leqslant$ 下的任何一个全序集.
\item 集合 $S$ 的所有子集的族，其偏序由包含关系确定（即：若 $A\subset B$，则 $A\preceq B$）.
\item $S$ 的一个子集族 $\mathcal{A}$，它关于有限交是封闭的，其偏序由反包含关系而确定（即，若 $A\supset B$，则 $A\preceq B$）.
\item 空间 $X$ 的所有闭子集的族，其偏序由包含关系确定.
\end{enumerate}
\item $J$ 的子集 $K$ 称为在 $J$ 中是\textbf{共尾的}（cofinal），若对于每一个 $\alpha\in J$，存在 $\beta\in K$，使得 $\alpha\preceq\beta$。证明若 $J$ 是一个有向集，并且 $K$ 在 $J$ 中是共尾的，则 $K$ 是一个有向集。
\item 设 $X$ 是一个拓扑空间，$X$ 中的一个\textbf{网}（net）是从一个有向集 $J$ 到 $X$ 的函数 $f$。如果 $\alpha\in J$，我们通常用 $x_{\alpha}$ 表示 $f(\alpha)$。而网 $f$ 本身则用符号 $(x_{\alpha})_{\alpha\in J}$ 表示，或者当指标集自明时，简单地记为 $(x_{\alpha})$。

称网 $(x_{\alpha})$ \textbf{收敛}（converge）于 $X$ 的点 $x$（写成 $x_{\alpha}\to x$），如果对于 $x$ 的每一个邻域 $U$，存在一个 $\alpha\in J$ 使得
\begin{align*}
\alpha\preceq\beta\Longrightarrow x_{\beta}\in U.
\end{align*}
证明：当 $J=\mathbb{Z}_{+}$ 时，这些定义就变成了我们所熟悉的那些相应概念。
\item 假设在 $X$ 和 $Y$ 中分别有
\begin{align*}
(x_{\alpha})_{\alpha\in J}\longrightarrow x\quad\text{和}\quad(y_{\alpha})_{\alpha\in J}\longrightarrow y.
\end{align*}
证明在 $X\times Y$ 中有 $(x_{\alpha}\times y_{\alpha})\to x\times y$。
\item 证明：若 $X$ 是 Hausdorff 的，则 $X$ 中的一个网最多收敛于一点。
\item \textbf{定理}\quad 设 $A\subset X$，则 $x\in\overline{A}$ 当且仅当存在 $A$ 中点的一个网收敛于 $x$。

[提示：为了证明蕴涵关系 $\Rightarrow$，取 $x$ 的所有邻域族为指标集，由反包含关系给出偏序。]
\item \textbf{定理}\quad 设 $f\colon X\to Y$。则 $f$ 是连续的当且仅当对于 $X$ 中任何一个收敛于 $x$ 的网 $(x_{\alpha})$，我们有网 $(f(x_{\alpha}))$ 收敛于 $f(x)$。
\item 设 $f\colon J\to X$ 是 $X$ 中的一个网，令 $f(\alpha)=x_{\alpha}$。若 $K$ 是有向集并且 $g\colon K\to J$ 是一个函数，满足条件：
\begin{enumerate}
\item $i\preceq j\Longrightarrow g(i)\preceq g(j)$,
\item $g(K)$ 在 $J$ 中共尾，
\end{enumerate}
则复合函数 $f\circ g\colon K\to X$ 称为 $(x_{\alpha})$ 的一个\textbf{子网}（subnet）。证明：若网 $(x_{\alpha})$ 收敛于 $x$，则它的任何一个子网也收敛于 $x$。
\item 设 $(x_{\alpha})_{\alpha\in J}$ 是 $X$ 中的一个网。称 $x$ 是网 $(x_{\alpha})$ 的一个\textbf{聚点}（accumulation point），如果对于 $x$ 的每一个邻域 $U$，满足条件 $x_{\alpha}\in U$ 的那些 $\alpha$ 的集合在 $J$ 中共尾。

\textbf{引理}\quad 网 $(x_{\alpha})$ 以 $x$ 为一个聚点当且仅当有 $(x_{\alpha})$ 的某一个子网收敛于 $x$。

[提示：为了证明蕴涵关系 $\Rightarrow$，令 $K$ 是所有偶对 $(\alpha,U)$ 的集合，其中 $\alpha\in J$，$U$ 是包含 $x_{\alpha}$ 的 $x$ 的一个邻域。当 $\alpha\preceq\beta$ 和 $V\subset U$ 时，定义 $(\alpha,U)\preceq(\beta,V)$，证明 $K$ 是有向集并且通过它定义一个子网。]
\item \textbf{定理}\quad $X$ 是紧致的当且仅当 $X$ 中的每一个网都有一个收敛子网。

[提示：为了证明蕴涵关系 $\Rightarrow$，令 $B_{\alpha}=\{x_{\beta}\mid\alpha\preceq\beta\}$。然后证明 $\{B_{\alpha}\}$ 满足有限交条件。为了证明蕴涵关系 $\Leftarrow$，令 $\mathcal{A}$ 为满足有限交条件的一个闭集族，$\mathcal{B}$ 为 $\mathcal{A}$ 中成员的所有有限交的族，并用反包含关系定义偏序。]
\item \textbf{推论}\quad 设 $G$ 是一个拓扑群，$A$ 和 $B$ 是 $G$ 的两个子集。若 $A$ 是 $G$ 的一个闭集，而 $B$ 是紧致的，则 $A\cdot B$ 在 $G$ 中是闭的。

[提示：假设 $G$ 是可度量化的，利用序列先给出一个证明。]
\item 验证：即使去掉有向集定义中的条件 (2)，上面这些习题的结论仍然成立。许多数学家就是在这种更一般的意义下使用``有向集''这个术语的。
\end{enumerate}